\section{Conclusion}\label{sec:conclusion}

We gave a parameterized account of MP-TSCFLP, and the picture that
emerges has two layers, a coverage layer, structural and strong, that
fires wherever a facility side sees a dimension discriminating it
through a $\{0,1\}$ incidence matrix, and a numeric layer, weak, that
dissolves under a unary encoding where no pair discriminates.
Table~\ref{tab:landscape} and Figure~\ref{fig:resultsmap} record the
resulting landscape entry by entry, from the size-parameter dichotomy
through the bracketing of the parameter $k$ to the kernelization and
width results, with witnesses on both sides of every border of the
binary-encoding map.

Five questions remain, and we regard them as the natural continuation of
this work. The first asks for a greedy $\ln|K|$ upper bound matching the
inapproximability of the coverage layer, the constructive half that would
close the approximability of that layer from above and below, with the
non-metric uncapacitated median problem \citep{Hochbaum1982} as the
precedent. A parameterized companion attaches to the same layer. The
reduction behind Theorem~\ref{thm:cover}, run with the enlarged
amplifier of Corollary~\ref{cor:coverinapprox}, maps the minimum cover
size onto the optimum value while shifting the parameter by one, so a
constant-factor \FPT\ approximation in $k$ for MP-TSCFLP would hand
back one for \textsc{Set Cover} parameterized by the cover size, which
is itself $\mathrm W[2]$-hard to approximate within any constant
\citep{LinRenSunWang2023}, and we conjecture that no such approximation
exists. The second concerns membership of the parameter $k$ in $\mathrm
W[2]$, since our sandwich leaves the exact class as the one open
classification of the project, and the barrier of binary arithmetic and
flow minimisation inside the verifier against weft two is of interest
beyond this model. The third would settle the unary status of $\{|J|,|L|\}$ and
of the two triples $\{|I|,|K|,|L|\}$ and $\{|J|,|K|,|L|\}$, the only gaps
of the sixteen-by-two dichotomy table, where the witnessing slices go
polynomial under unary data and no other hard slice is yet known. The
fourth is whether a polynomial kernel exists in the full parameter
$|I|+|J|+|K|+|L|$ with general numbers, the exact point where the
yes-certificate comparison turns bilinear and both weight families would
have to be reduced at once, a genuine methodological frontier of weighted
kernelization. The fifth turns to transport costs structured
beyond separability, uniform per side, metric or of low rank, fall on the
tractable or the hard side, refining the frontier over the cost classes
that appear in real data. Every hardness construction in the paper uses
transport costs without metric structure, so the map may bind
differently on distance-based data, and the fifth open problem asks
exactly that question. Progress on any one of the five would redraw a
border of the landscape assembled here.

For practice the reading is brief and concrete. The a priori bound
$k^\ast$ is computed by a sort and a prefix scan of the capacities per
product and side, and by nothing costlier,
lower-bounds the
facilities any feasible design must open, and is the very threshold on
which our enumeration prunes. The completed campaign
tests it against solver effort on the full benchmark and finds a
direction a naive reading would not guess, for larger bounds accompany
smaller terminal gaps, in an association that is negative in every size
group and survives multiplicity correction pooled and in the
five-product groups. The log analysis places that association at the
root relaxation rather than at the search, since almost every run ends
at the time limit with its bound close to the root value, so on this
benchmark the bound is best read as an annotation of relaxation
tightness and not as a forecast of search effort.
